\section{Application Architecture}
\label{sec:apparch}

\subsection{Domain Model}

The domain is built on value types so that state is trivially snapshot-able,
diffable, and thread-safe.

\begin{lstlisting}[language=Swift, caption={Core domain types.}, label={lst:domain}]
/// A complete, self-contained description of how to develop one image.
struct Recipe: Codable, Hashable, Sendable {
    var negativeStock : StockRef
    var printStock    : PrintStockRef
    var chemistry     : ChemistryRef

    var capture   : CaptureStage      // exposure comp, WB, rating
    var develop   : DevelopStage      // push/pull, time, temp, agitation
    var interlayer: InterlayerStage
    var printing  : PrintStage        // printer lights, density, sat.
    var grain     : GrainStage
    var halation  : HalationStage
    var optical   : OpticalStage
    var aging     : AgingStage
    var output    : OutputStage       // color space, surround, crop

    var format    : FilmFormat        // drives all physical scaling
    var seed      : UInt32            // grain / dust / leak realization
}

/// Stable identity for every scalar the pipeline consumes.
enum ParameterKey: String, CaseIterable, Codable, Sendable {
    case exposureCompensation, filmSpeed, whiteBalanceTemp, whiteBalanceTint
    case pushPull, developmentTime, developmentTemp, agitation
    case couplerActivity
    case printerLightR, printerLightG, printerLightB, printDensity
    case saturationDensity, highlightRolloff, shadowLift
    case neutralAxisWarm, neutralAxisTint, silverRetention
    case grainAmount, grainSize, grainSeed
    case halationThreshold, halationRadius, halationIntensity, halationBias
    case diffusionAmount, vignetteAmount, chromaticAberration
    case distortionK1, distortionK2
    case dustDensity, scratchDensity, edgeFog, lightLeak, expiredAge
    // ...
}
\end{lstlisting}

\texttt{Recipe} is \texttt{Codable}, which gives sidecar export
(FR-17) and database persistence for free, and \texttt{Hashable}, which gives
cheap change detection. Every field has a documented default, and a
\texttt{Recipe} constructed with defaults plus a stock reference is a valid,
renderable state --- there is no partially-initialized recipe.

\subsection{Parameter Resolution}

A rendered image's parameters are the result of layering four sources, resolved
in order:

\begin{enumerate}
\item Stock profile defaults (from the profile bundle).
\item Chemistry profile modulation (Section~\ref{sec:development}).
\item Recipe values (the user's edit).
\item Transient overrides (preview quality reductions, gesture state).
\end{enumerate}

Resolution produces a \texttt{ResolvedParameters} struct --- a flat, contiguous
value with no optionals --- which is what the render graph consumes. The
separation matters: the UI edits a sparse \texttt{Recipe} where "unset" is
meaningful (so that changing stock changes the defaults for untouched
parameters), while the renderer consumes a dense resolved value where everything
has a number.

\subsection{Edit History}

History is implemented as a list of \texttt{Recipe} snapshots rather than a
command stack. A \texttt{Recipe} is a few hundred bytes; 200 snapshots is under
100\,KB, which is cheaper than the bookkeeping a command pattern requires and is
immune to the classic inverse-command bug where undo does not exactly restore
prior state.

Snapshots are \emph{coalesced} by a policy that merges consecutive edits to the
same parameter key within a 700\,ms window, so that a slider drag produces one
undo step rather than four hundred. The coalescing rule is expressed on the key,
not on time alone --- switching parameters always commits a step, so the user's
mental model of undo matches the visual grouping of controls.

\begin{lstlisting}[language=Swift, caption={History with key-aware coalescing.}]
struct History {
    private(set) var stack: [Entry] = []
    private(set) var cursor: Int = -1
    private var lastKey: ParameterKey?
    private var lastAt: ContinuousClock.Instant?

    mutating func record(_ r: Recipe, key: ParameterKey?, now: ContinuousClock.Instant) {
        let coalesce = key != nil && key == lastKey
            && lastAt.map { now - $0 < .milliseconds(700) } ?? false
        if coalesce, cursor >= 0 {
            stack[cursor].recipe = r          // replace in place
        } else {
            stack.removeSubrange((cursor + 1)...)   // drop redo branch
            stack.append(Entry(recipe: r, key: key, at: now))
            cursor = stack.count - 1
            if stack.count > 200 { stack.removeFirst(); cursor -= 1 }
        }
        lastKey = key; lastAt = now
    }
}
\end{lstlisting}

\subsection{The Render Actor State Machine}

The coalescing behavior described in Section~\ref{sec:arch} is specified here
precisely, because it is the difference between a responsive and a laggy editor
and it is easy to get subtly wrong.

\begin{lstlisting}[language=Swift, caption={Render actor with single-slot coalescing.}, label={lst:renderactor}]
actor RenderCoordinator {
    private var inFlight: Bool = false
    private var pending: RenderRequest?
    private let graph: RenderGraph

    /// Called from the main actor on every parameter change. Never blocks.
    func submit(_ request: RenderRequest) {
        pending = request                 // overwrite: only latest matters
        guard !inFlight else { return }   // completion handler will pick it up
        drain()
    }

    private func drain() {
        guard let req = pending else { inFlight = false; return }
        pending = nil
        inFlight = true

        let cb = queue.makeCommandBuffer()!
        do { try graph.encode(req, into: cb) }
        catch { inFlight = false; report(error); return }

        cb.addCompletedHandler { [weak self] _ in
            Task { await self?.drain() }   // pick up whatever arrived meanwhile
        }
        cb.commit()
    }
}
\end{lstlisting}

Three properties follow and are asserted by test:

\begin{itemize}
\item At most one command buffer is in flight, bounding memory and latency.
\item Intermediate parameter values during a fast drag are dropped rather than
queued, so latency does not accumulate.
\item The final value of a drag is always rendered, because the last
\texttt{submit} sets \texttt{pending} and the completion handler will drain it.
The naive "skip if busy" implementation loses the final value and leaves a stale
image, which is a defect that only appears when the user flicks a slider and
lets go --- and which is therefore routinely shipped.
\end{itemize}

\subsection{Capture Subsystem}

Capture uses \texttt{AVCaptureSession} with a \texttt{AVCapturePhotoOutput}
configured for Bayer or ProRAW output depending on device capability.

\begin{lstlisting}[language=Swift, caption={ProRAW capture configuration.}]
func makeSettings(output: AVCapturePhotoOutput) -> AVCapturePhotoSettings {
    // Prefer Apple ProRAW (linear DNG with Apple's computational stack applied
    // to the mosaic) over plain Bayer RAW: it carries the multi-frame
    // noise reduction we do not wish to reimplement, while remaining linear.
    if output.isAppleProRAWSupported {
        output.isAppleProRAWEnabled = true
    }
    let fmt = output.availableRawPhotoPixelFormatTypes.first {
        AVCapturePhotoOutput.isAppleProRAWPixelFormat($0)
    } ?? output.availableRawPhotoPixelFormatTypes.first!

    let s = AVCapturePhotoSettings(rawPixelFormatType: fmt,
                                   processedFormat: [AVVideoCodecKey: AVVideoCodecType.hevc])
    s.photoQualityPrioritization = .quality
    s.isHighResolutionPhotoEnabled = true
    s.embeddedThumbnailPhotoFormat = [AVVideoCodecKey: AVVideoCodecType.jpeg]
    return s
}
\end{lstlisting}

Manual controls are exposed through \texttt{AVCaptureDevice} lock:
\texttt{setExposureModeCustom} for shutter and ISO,
\texttt{setFocusModeLocked} for manual focus with a lens-position slider, and
\texttt{setWhiteBalanceModeLocked} with device gains computed from a temperature
and tint pair. The camera preview shows a live approximation of the selected
recipe (Section~\ref{sec:liveview}), not the raw sensor image, since seeing the
intended result while composing is the primary reason a photographer would use
this application's camera rather than the system one.

\subsection{Live Preview Approximation}
\label{sec:liveview}

The full pipeline at video rates on the preview stream is achievable but is not
a good use of thermal budget while composing. The live view therefore runs a
reduced graph:

\begin{itemize}
\item The pointwise chain runs at full fidelity, via the same baked 3D LUT
(Section~\ref{sec:lutbake}) applied to the preview stream. This is a single
texture fetch and is what determines the color, which is what the user is
judging.
\item Halation runs with a two-level pyramid instead of four.
\item Grain and aging are disabled, with a small badge indicating so.
\item Optical simulation runs vignetting only.
\end{itemize}

Measured cost is 2.8\,ms per frame at $1920\times1440$ preview resolution on an
A17 Pro, comfortably inside the 33\,ms frame budget with substantial thermal
headroom.

\subsection{Import Pipeline}

Imported files traverse: security-scoped URL resolution; format probe;
\texttt{CIRAWFilter} construction and metadata extraction; copy into the
application's managed store (or, for Photos assets, retention of a
\texttt{PHAsset} local identifier plus a cached decode); thumbnail generation at
three sizes; and a database insert. The entire operation runs on the I/O queue
and reports progress; a batch import of 200 DNGs completes in approximately
40\,s, dominated by thumbnail decode.

Imported files whose color characterization cannot be determined (an unusual
TIFF, a DNG from an unknown camera) are flagged and given a conservative
assumed-sRGB interpretation, with a visible advisory. Silently guessing is worse
than saying so.

\subsection{Error Handling Policy}

Errors are classified and handled by class rather than individually:

\begin{table}[htbp]
\caption{Error Classification}
\label{tab:errors}
\begin{center}
\renewcommand{\arraystretch}{1.15}
\begin{tabularx}{\columnwidth}{@{}l>{\raggedright\arraybackslash}X@{}}
\toprule
\textbf{Class} & \textbf{Policy} \\
\midrule
Programmer error & Trap in debug, log and degrade in release. Graph cycles, missing pipeline states, impossible parameter combinations. \\
Resource exhaustion & Reduce quality tier and retry once; then surface. Applies to heap allocation and export tiling. \\
Device loss / GPU fault & Recreate device objects, rebuild caches, re-render. Report once per session. \\
User data error & Surface with a specific, actionable message. Unsupported file, corrupt DNG, missing permission. \\
Transient & Retry with backoff, silent. Photos library access contention, file coordination. \\
\bottomrule
\end{tabularx}
\end{center}
\end{table}
